\section{Polynomial Linear Regression}

Given data-points \((x_i, y_i) \in \R^2, \ i = 1, \dots, m\) we define the \emph{mean square error} as

\[
    \tilde{e} = \sqrt{ \sum_{i = 1}^{m} {(f(x_i) - y_i)}^2 } = \sqrt{\sum_{i = 1}^{m} e_i^2} = \metric{e}_2, \quad e = 
    \begin{bmatrix}
        e_1 \\ 
        \vdots \\ 
        e_m
    \end{bmatrix}   
\]

Now instead of using a linear term we define a polynomial of degree \(n- 1\): 

\[
    f(x) = p_1 + p_2 x + \cdots + p_m x^{n - 1}.
\]

Hence we wan to minimize the expression 

\[
    \metric{e}_2 = \metric{
        \begin{bmatrix} 
            f(x_1) - y_1 \\ 
            \vdots \\ 
            f(x_m) - y_m
        \end{bmatrix}}_2
\]

Which delivers the following system 

\[
    \metric{
        \begin{bmatrix} 
            1 & x_1 & \cdots & x_1^{n - 1} \\ 
            \vdots  & \vdots & & \vdots \\ 
            1 & x_m & \cdots & x_m^{n -1}
        \end{bmatrix}
        \begin{bmatrix}
            p_1 \\ 
            \vdots \\ 
            p_n
        \end{bmatrix} 
        -
        \begin{bmatrix}
            y_1 \\ 
            \vdots \\ 
            y_m
        \end{bmatrix}
        }_2 = 
        \metric{Ap - y}_2
\]

Thus, given \(A \in \R^{m \times n}\) such that \(\rank(A) = n\) and \(m \ge n\), we have the minimization problem 

\[
        \min_{x \in \R^n} \metric{Ap - y}_2
\]

whose solution \(\tilde{x} \in \R^n\) is also the solution of the \emph{normal equation}

\[
        A^T A \tilde{x} = A^T b 
\]

whereas \(A^T A \in \R^{n \times n}\) is regular.

\textbf{Proof:}

We want to first show that \(A^T A\) is regular hence, we do the following: 

If \(A^T Ax = 0\) then 


\[
    0 = x^T A^T Ax = {(Ax)}^T Ax = \metric{Ax}_2^2 \implies Ax = 0
\]

Because \(A\) has maximal rank \(x = 0\), hence \(A^T A\) is regular.

Now we can show that \(\tilde{x} = (A^T A)^{-1}A^T b\) minimizes \(\metric{Ax - b}_2^2\)

\begin{align*}
    J(x) = \metric{Ax - b}_2 &= (Ax - b)^T(Ax - b) \\
     &= ((Ax)^T - b^T)(Ax - b) \\
     &= x^T(A^TAx - Ab^T) + b^T(b - Ax)
\end{align*}

Hence, \(J(\tilde{x}) = \tilde{x}^T((A^T)(A^TA)^{-1} A^Tb - A^Tb) + b^T(b - A\tilde{x}) = b^T(b - A\tilde{x})\).

If we assume there exits another solution \(\tilde{y} =  \tilde{x} ++ v, \ v \in \R^n, \ v\ne 0\) then 

\begin{align*}
    J(\tilde{y}) &= (\tilde{x} + v)^T(A^T A(\tilde{x} + v) - A^Tb) + b^T (b - A(\tilde{x} + v)) \\ 
     &= (\tilde{x} + v)^T(A^T A\tilde{x} + A^TAv - A^Tb) + b^T (b - A\tilde{x} + Av)) \\ 
     &= (\tilde{x} + v)^T A^T A\tilde{x} + A^TAv  + b^T (b - A\tilde{x}) - b^T A^ v \\ 
     &=  v^TA^T A(\tilde{x} + v)^T - v^TA^Tb  + b^T (b - A\tilde{x}) \\ 
     &=  v^T(A^T A\tilde{x}(\tilde{x}  - A^Tb)^T + v^TA^Tb  + b^T (b - A\tilde{x}) \\ 
     &= \metric{Av}_2^2 + J(\tilde{x})
\end{align*}

Because \(v \ne 0\) and \(A\) has the maximal rank, \(Av \ne 0\) and \(\metric{Av}_2^2 > 0\) hence 

\[
    J(\tilde{y}) = \metric{A(\tilde{y} - \tilde{x})}_2^2 + J(\tilde{x}) > J(\tilde{x})
\]

Therefore, \(\tilde{x}\) is the solution with the minimal distance.

\QED

\textbf{Example:}

Given the points \((0, 1), (2, 2), (3, 2)\) find the minimal solution to the minimization problem using \(f(x) = p_1 + p_2 x, \,(n = 2)\).
Notice that our basis is \(1, \, x\); hence:

\[
        A = 
        \begin{bmatrix}
            1 & x_1 \\ 
            1 & x_2 \\ 
            1 & x_3
        \end{bmatrix}
        = 
        \begin{bmatrix}
            1 & 0 \\ 
            1 & 2 \\ 
            1 & 3
        \end{bmatrix},
        \quad 
        b = 
        \begin{bmatrix}
            y_1 \\ 
            y_2 \\ 
            y_3
        \end{bmatrix}
        =
        \begin{bmatrix}
            1 \\ 
            2 \\ 
            2
        \end{bmatrix}, 
        \quad 
        x = 
        \begin{bmatrix}
            p_1 \\ 
            p_2
        \end{bmatrix}
\]

Hence for \(A^T A x = A^T b\) we get 

\[
        A^TA = 
        \begin{bmatrix}
            1 & 1 & 1 \\ 
            0 & 2 & 3
        \end{bmatrix}
        \begin{bmatrix}
            1 & 0 \\ 
            1 & 2 \\ 
            1 & 3
        \end{bmatrix}
        =  
        \begin{bmatrix}
            3 & 5 \\ 
            5 & 13
        \end{bmatrix}
\]

and 

\[
        A^T b
        =
        \begin{bmatrix}
            1 & 1 & 1 \\ 
            0 & 2 & 3        
        \end{bmatrix}
        \begin{bmatrix}
            y_1 \\ 
            y_2 \\ 
            y_3
        \end{bmatrix}
        = 
        \begin{bmatrix}
            5 \\ 
            10
        \end{bmatrix}
\]

Hence by solving 

\[
        \begin{bmatrix}
            3 & 5 \\ 
            5 & 13
        \end{bmatrix}
         \begin{bmatrix}
            p_1 \\ 
            p_2
        \end{bmatrix}
        = 
        \begin{bmatrix}
            5 \\ 
            10
        \end{bmatrix}
\]

using Gaussian elimination we get \(p_1 = \frac{15}{14}\) and \(p_2 = \frac{5}{14}\); giving us the final polynomial 

\[
        f(x) = \frac{5}{14} + \frac{15}{14}x
\]


\section{Pseudo Inverse}

Given \(A \in \R^{m \times n}\) and \(b \in \R^m\), there is exactly one solution to the minimization problem \(x\ in \R^n\) which is dependent on the 
matrix \(A^+ \in \R^{n \times m}\) with the property 

\[
        A^+ b = x, \ \forall b \in \R^m.
\]

This matrix is the \emph{pseudo inverse of  A} and obeys the following properties 

\begin{itemize}

    \item \((A^+ A)^T = A^+ A\)
   
    \item \((AA^+)^T = AA^+ \)

    \item \(A^+ A A^+ = A^+\)
    
    \item \(A A^+ A = A\)

\end{itemize}

\section{Singular Value Decomposition}

Given \(A \in \R^{m \times n}\), then there exists matrices \(U \in \R^{m \times m}, \ V \in \R^{n \times n}\) such that 

\[
    A = U \Sigma V^T
\]

where

\[
    \Sigma \in \R^{m \times n}, \ \Sigma = \mathrm{diag}(\sigma_1, \dots, \sigma_p), \ p = \min(m,n)
\]

and 

\[
    \sigma_1 \ge \sigma_2 \ge \dots \ge \sigma_p \ge 0.
\]

This is the \emph{singular value decomposition of \(A\)}. 

For the case \(A \in \R^{n \times n}\) this is just a diagonalisation.

\subsubsection*{Properties Of The SVD}

\begin{itemize}

    \item \(V\Sigma^T U^T\) is SVD of \(A^T\) and has the same \emph{singular values} as \(A\).

    \item If \(\sigma_1 \ge \dots \ge \sigma_r > \sigma_{r + 1} = \dots  = \sigma_p = 0\), then \(\rank(A) = r\).

    \item \(\sigma_1^2, \dots \sigma_p^2\) are the eigenvalues of \(AA^T, \ A^TA\).

    \item The columns of \(U,\ V\) are the eigenvectors of \(AA^T, \ A^TA\).

    \item If \(r = \rank(A), \ A = U \Sigma V^T\) with \(\Sigma = \mathrm{diag}(\sigma_1, \dots, \sigma_r, 0, \dots, 0) \in \R^{m \times n}\), then
    
    \[
        A^+ = V \Sigma^+ U^T, \quad \Sigma^+ = \mathrm{diag}(\frac{1}{\sigma_1}, \dots, \frac{1}{\sigma_r}, 0, \dots, 0,) \in \R^{n \times m}
    \]

\end{itemize}

Hence, to compute the singular value decomposition we need to compute the eigenvalues and eigenvectors of \(AA^T\) or \(A^TA\) and take the 
square root of its eigenvalues to get the \emph{singular values}.

\section{SVD Regularization}

Given \(A \in \R^{m \times n}\) with a SVD \(A = U \Sigma V^T\) with 

\[
    \Sigma = \mathrm{diag}(\frac{1}{\sigma_1}, \dots, \frac{1}{\sigma_r}, 0, \dots, 0,) \in \R^{m \times n}.
\]

For \(x = A^+b, \, \Delta x = A^+ \Delta b\), the following approximations are valid 

\[
    \frac{\metric{\Delta x}_2}{\metric{x}_2} \le \frac{\sigma_1}{\sigma_r} \frac{\metric{\Delta b}_2}{\metric{P_r U^T b}_2},
\]

with \(P_r = \mathrm{diag}(1_1, \dots, 1_r, 0, \dots, 0) \in \R^{m \times m}\). There is always a \(b\) and \(\Delta b\) such that

\[
    \frac{\metric{\Delta x}_2}{\metric{x}_2} = \frac{\sigma_1}{\sigma_r} \frac{\metric{\Delta b}_2}{\metric{b}_2}
\]

We also define the \emph{condition} of our matrix as \(\kappa_2(A) = \frac{\sigma_1}{\sigma_r}\) if all of our singular values are greater than zero.

\subsection{Truncated SVD (TSVD)}

Given \(A \in \R^{n \times n}\) with a SVD \(A = U \Sigma V^T\), 

\[
    \Sigma = \mathrm{diag}(\sigma_1, \dots, \sigma_r, 0, \dots, 0,) \in \R^{m \times n}, \ \sigma_1 \ge \dots \ge \sigma_r > 0.
\]

For \(\alpha > 0\) is \(A_{\alpha} = U \Sigma_{\alpha} V^T\)

\[
    \Sigma_{\alpha}= \mathrm{diag}(\sigma_1, \dots, \sigma_k, 0, \dots, 0,) \in \R^{m \times n}, \ \frac{\sigma_k}{\sigma_1} \ge \sqrt{\alpha} \ge \frac{\sigma_{k + 1}}{\sigma_1}.
\]

is the \emph{truncated singular value decomposition} of \(A\).

From this definition we also get the fact \(\kappa_2 (A_{\alpha}) \le \frac{1}{\sqrt{\alpha}}\).

\section{Tikhonov Regularization}

When computing \(x = A^+ b\) we get the equivalent minimization problem 

\[
    \min \metric{Ax - b}_2, \quad \min \metric{x}.
\]

By combining then in a expression we get 

\begin{align*}
    J_{\alpha}(x) &= \metric{Ax - b}_2^2 + \alpha \metric{x}_2^2 \\ 
    &= \metric{
        \begin{bmatrix} 
            Ax \\ 
            \sqrt{\alpha}x
        \end{bmatrix}
        - 
        \begin{bmatrix}
            b \\ 
            0
        \end{bmatrix}
        }_2^2 \\ 
    &= \metric{
        \begin{bmatrix} 
            A \\ 
            \sqrt{\alpha}I
        \end{bmatrix}x
        - 
        \begin{bmatrix}
            b \\ 
            0
        \end{bmatrix}
        }_2^2 \\ 
    &= \metric{A_{\alpha}x - b_{\alpha}}_2^2
\end{align*}

The rank is preserved by \(\alpha > 0\), hence we can minimize the normal equation

\[
    A_{\alpha}^T A_{\alpha} x_{\alpha} = A_{\alpha}^T b_{\alpha}
\]

We can rewrite as

\[
    [A^T, \alpha I]
    \begin{bmatrix}
        A \\ 
        \alpha I
    \end{bmatrix}
    x_{\alpha}
    = [A^T, \alpha I] b_{\alpha}
\]

giving us 

\[
    x_{\alpha} = (A^T A + \alpha I)^{-1} A^T b
\]

Formally, for \(A \in \R^{m \times n},\ \alpha > 0\) is 

\[
    x_{\alpha} = (A^T A + \alpha I)^{-1} A^T b
\]

the \emph{Tikhonov regularization} of \(x = A^+ b\) with parameter \(\alpha\).








